% Start preamble
\documentclass[12pt,a4paper]{article}
\usepackage{geometry}
 \geometry{
 a4paper,
 total={170mm,257mm},
 left=20mm,
 top=20mm,
 }
\usepackage[utf8]{inputenc}
\usepackage[T1]{fontenc}
\usepackage[pdftex]{graphicx}
\graphicspath{{./}}
\usepackage{enumitem}
\usepackage{pdfpages}
\usepackage{hyperref}
\usepackage{tikz}
\usepackage{attachfile}
\usepackage{epstopdf}
\usepackage{array}
\usepackage{multirow}
\usepackage{multicol}
%\usepackage{makecell}
%\usepackage{multirow}
\usepackage[table]{xcolor}

\usepackage{float}
%\usepackage[table]{xcolor,colorbl}
\setlength{\textwidth}{17cm}
\setlength{\oddsidemargin}{-0.5cm}
\setlength{\evensidemargin}{-0.5cm}
%\setlenght{\headsep}{0cm}
\setlength\parindent{0pt}
%\setlength{\extrarowheight}{3pt}
\usepackage{listings}
%\usepackage{xcolor}

\input{arduinoLanguage.tex}
%%%%%% Counting oppgaves %%%%%%
 \newcount\questnum \questnum=0
 \def\oppgave{
            \advance\questnum by 1
	    \ifthenelse{\questnum>0\AND \questnum<9}
	    {
                \vskip 1cm
		\textbf{Oppgave}\hskip 5pt\the\questnum \hfill \hfill(6p)
		\vskip 3pt
		\hrule
	\vskip 0.5cm}
	{
                \vskip 1cm
		\textbf{Oppgave}\hskip 5pt \the\questnum \hfill \hfill(12p)
		\vskip 3pt \hrule \vskip 0.5cm }

		}

\begin{document}
\title{Konte Prøve for 2INTA\\ Industrielektronikk 02\vskip 1cm Navn:\leaders\vrule height 0pt depth 1pt \hfill \null}
\author{Faglærer: Fred-Olav Mosdal 90507684\\Hjelpemidler: Kalkulator og verkstedhåndbok}
\date{19.11.2024}
\maketitle
	$$\includegraphics[height=15cm]{../output/nogpl/elektro_tr.png}$$

\newpage
\oppgave{}%1
\vskip 2.5pt 
Hva skjer i en krets bestående av en spenningskilde og en resistans når:

\begin{enumerate}
\item Spenningen tredobles
\vskip 5pt 
\begin{tikzpicture}
	\draw[step=0.5cm,gray!20,very thin]  grid (16,1) ;
\end{tikzpicture}
\vskip 2.5pt 
\item Spenningen reduseres med 75\%
\vskip 5pt 
\begin{tikzpicture}
	\draw[step=0.5cm,gray!20,very thin]  grid (16,1) ;
\end{tikzpicture}
\vskip 2.5pt 
\item Resistansen dobbles
\vskip 5pt 
\begin{tikzpicture}
	\draw[step=0.5cm,gray!20,very thin]  grid (16,1) ;
\end{tikzpicture}
\vskip 2.5pt 
\item Resistansen reduseres med 35\%
\vskip 5pt 
\begin{tikzpicture}
	\draw[step=0.5cm,gray!20,very thin]  grid (16,1) ;
\end{tikzpicture}
\vskip 2.5pt 
\item Spenningen dobbles og resistansen dobbles
\end{enumerate}

\oppgave{}%2
\vskip 2.5pt 

Hva oppnår vi ved å seriekoble to batterier?
\vskip 5pt 
\begin{tikzpicture}
	\draw[step=0.5cm,gray!90,very thin]  grid (17,4) ;
\end{tikzpicture}
\vskip 2.5pt 
\oppgave{}%3
\vskip 2.5pt 
\begin{enumerate}
\item Hva er benevnelsen og enheten til spenning?
\\
\begin{tikzpicture}
	\draw[step=0.5cm,gray!90,very thin]  grid (16,2);
\end{tikzpicture}
\item Hva er benevnelsen og enheten til strøm?
\\
\begin{tikzpicture}
	\draw[step=0.5cm,gray!90,very thin]  grid (16,2);
\end{tikzpicture}
\item Hva er benevnelsen og enheten til resistans?
\\
\begin{tikzpicture}
	\draw[step=0.5cm,gray!90,very thin]  grid (16,2);
\end{tikzpicture}
\end{enumerate}
\vskip 2.5pt 
\newpage
\oppgave{}%5
\vskip 2.5pt 
Hvor stor strøm går det gjennom en resistans på $1000\Omega$ når det ligger en spenning på 220V over den? 
\vskip 5pt 
\begin{tikzpicture}
	\draw[step=0.5cm,gray!90,very thin]  grid (17,3) ;
\end{tikzpicture}
\vskip 2.5pt 
\oppgave{}%5
\vskip 2.5pt 
Hvor stor spenning kreves for å drive 5A gjennom en resistans på $200\Omega$? 
\vskip 5pt 
\begin{tikzpicture}
	\draw[step=0.5cm,gray!90,very thin]  grid (17,3) ;
\end{tikzpicture}
\vskip 2.5pt 
\oppgave{}%4
\vskip 2.5pt 
Hvor stor resistans er det i en panelovn om den er tilkoblet spenning på U=10V og det går en strøm på I=5A i den?
\vskip 5pt 
\begin{tikzpicture}
	\draw[step=0.5cm,gray!90,very thin]  grid (17,4) ;
\end{tikzpicture}
\vskip 2.5pt 
\newpage
\oppgave{}%6
\vskip 2.5pt 
Vi har en seriekobling med $R_{1}=2200\Omega$ og $R_{2}$. Når vi kobler 24V til seriekoblingens ytterpunkter, går
det 5mA i kretsen. 

\begin{enumerate}
\item Tegn skjema for kretsen
\vskip 2.5pt 
\begin{tikzpicture}
	\draw[step=0.5cm,gray!90,very thin]  grid (15,8) ;
\end{tikzpicture}
\item Hvor stor strøm går det gjennom $R_{1}$ og $R_{2}$ 
\vskip 2.5pt 
\begin{tikzpicture}
	\draw[step=0.5cm,gray!90,very thin]  grid (15,2) ;
\end{tikzpicture}
\item Hvor stor er kretsens totale resistans?
\vskip 2.5pt 
\begin{tikzpicture}
	\draw[step=0.5cm,gray!90,very thin]  grid (15,2) ;
\end{tikzpicture}
\item Hvor stor er $R_{2}$?
\vskip 2.5pt 
\begin{tikzpicture}
	\draw[step=0.5cm,gray!90,very thin]  grid (15,2) ;
\end{tikzpicture}
\item Hvor stor spenning ligger det over hver av motstandene?
\vskip 2.5pt 
\begin{tikzpicture}
	\draw[step=0.5cm,gray!90,very thin]  grid (15,2) ;
\end{tikzpicture}
\end{enumerate}
\vskip 5pt 
\vskip 2.5pt 
\newpage
\oppgave{}%7
\vskip 2.5pt 
To resistanser er parallellkoplet $R_1=500\Omega$ og $R_2=750\Omega$. Resistansene er tilkoplet en spenning på 4V. 
\vskip 5pt 
\begin{enumerate}
\item Tegn skjema for kretsen
\vskip 2.5pt 
\begin{tikzpicture}
	\draw[step=0.5cm,gray!90,very thin]  grid (15,8) ;
\end{tikzpicture}
\item Hva blir kretsens totale resistans?
\vskip 2.5pt 
\begin{tikzpicture}
	\draw[step=0.5cm,gray!90,very thin]  grid (15,3) ;
\end{tikzpicture}
\item Hva blir kretsens hovedstrøm?
\vskip 2.5pt 
\begin{tikzpicture}
	\draw[step=0.5cm,gray!90,very thin]  grid (15,2) ;
\end{tikzpicture}
\item Hvor stor er strømmen i $R_{1}$?
\vskip 2.5pt 
\begin{tikzpicture}
	\draw[step=0.5cm,gray!90,very thin]  grid (15,2) ;
\end{tikzpicture}
\item Hvor stor er strømmen i $R_{2}$?
\vskip 2.5pt 
\begin{tikzpicture}
	\draw[step=0.5cm,gray!90,very thin]  grid (15,2) ;
\end{tikzpicture}
\end{enumerate}
\vskip 2.5pt 
\newpage
\oppgave{}%9
\vskip 2.5pt 
Praktisk oppgave \\
I denne oppgavn er $R_1=2200\Omega$, $R_2=2700\Omega$ og $R_3=3300\Omega$
	\Large
\begin{center}
	\begin{tabular}{| m{10.2cm} |m{1.5cm} |m{1.5cm} |m{1.5cm} |} 
\hline
	\multicolumn{4}{|c|}{\textbf{\cellcolor[HTML]{D5D5D5}Måle- og koblingsoppgaver}} \\
\hline
\hline
\rowcolor [HTML]{D5D5D5}
Oppgave	&Ikke bestått&OK&Bestått\\ \hline

%	\multicolumn{4}{|c|}{\textbf{\cellcolor[HTML]{D5D5D5}Enkel motstand}} \\
%\hline
		%Koble $R_1$ til spenning&&&\\ \cline{1-4}
		%Mål spenningen over motstanden&&&\\ \cline{1-4}
		%Mål strømmen igjennom motstanden&&&\\ \cline{1-4}
	\multicolumn{4}{|c|}{\textbf{\cellcolor[HTML]{D5D5D5}Seriekobling}} \\
\hline
		Koble $R_1$ og $R_2$ i serie&&&\\ \cline{1-4}
		Mål spennigen over $R_1$&&&\\ \cline{1-4}
		Mål spennigen over $R_2$&&&\\ \cline{1-4}
		Mål strømmen i seriekoblingen &&&\\ \cline{1-4}
		Mål den totale resistansen i kretsen&&&\\ \cline{1-4}
	\multicolumn{4}{|c|}{\textbf{\cellcolor[HTML]{D5D5D5}Parallellkobling}} \\
		Koble $R_2$ og $R_3$ i parallell&&&\\ \cline{1-4}
		Mål strømmen i parallellkoblingen &&&\\ \cline{1-4}
		Mål strømmen i $R_2$&&&\\ \cline{1-4}
		Mål strømmen i $R_3$&&&\\ \cline{1-4}
		Mål den totale resistansen i kretsen&&&\\ \cline{1-4}
\end{tabular}
\end{center}
\vskip 2.5pt 
\newpage
\end{document}
